\documentclass{article}
\usepackage[T1]{fontenc}
\usepackage{lmodern}
\usepackage{graphicx}
\usepackage{amsmath,amssymb}
\usepackage[a4paper,margin=2cm]{geometry}
\usepackage[absolute,overlay]{textpos}
\setlength{\TPHorizModule}{1mm}
\setlength{\TPVertModule}{1mm}
\usepackage[indonesian]{babel}
\usepackage{hyperref}
\setlength{\emergencystretch}{2em}
% unit: Halaman 82

\begin{document}

% === Halaman 82 (llm) ===

\section*{2. Program Steganografi Audio dengan Python}

\begin{itemize}
    \item Kode program Python berikut ini menyisipkan pesan dengan metode LSB ke dalam file audio
\end{itemize}

\textbf{A. Program penyisipan pesan}
\begin{itemize}
    \item Input:
    \begin{enumerate}
        \item File audio tidak terkompresi (format WAV)
        \item Pesan (diketik langsung dari keyboard)
    \end{enumerate}
    \item Output: stego audio (format WAV)
\end{itemize}

\textbf{B. Program ekstraksi pesan}
\begin{itemize}
    \item Input: Stego audio (format WAV)
    \item Output: pesan
\end{itemize}

\end{document}
